\documentclass[11pt]{article}
% \documentclass{assignment}

\usepackage[margin=1in]{geometry} 
\usepackage{amsmath,amsthm,amssymb, graphicx, multicol, array}
\usepackage{ctex} 
\usepackage{fancyhdr}
\usepackage{fontspec}
\usepackage{fourier-orns} % 
 
\newcommand{\N}{\mathbb{N}}
\newcommand{\Z}{\mathbb{Z}}
 
\theoremstyle{remark}
\newtheorem{problem}{}


% \newenvironment{problem}[2][]{\begin{trivlist}
% \item[\hskip \labelsep {\bfseries #1}\hskip \labelsep {\bfseries #2}]}{\end{trivlist}}

\renewcommand{\headrule}{%
\vspace{-8pt}\hrulefill
\raisebox{-2.0pt}{\quad\decofourleft\decotwo\decofourright\quad}\hrulefill}


\newcommand{\duedate}{\zhdate{2026/12/01}}
\pagestyle{fancy}

% \fancyhf{}
% \fancyhead[LE]{\nouppercase{\rightmark\hfill\leftmark}}
% \fancyhead[RO]{\nouppercase{\leftmark\hfill\rightmark}}

%... then configure it.
\fancyhead{} % clear all header fields
% \fancyhead[LE]{\textbf{PHYS}}
% \fancyhead[L]{\textbf{\leftmark}}
% \fancyhead[L]{\chaptermark{}}
% \fancyhead[L]{\sectionmark{}}
% \fancyhead[LR]{\textbf{\leftmark}}
\fancyhead[L]{数学物理方法}
\zhnumsetup{style={Traditional,Financial}}
\fancyhead[R]{课后习题捌} 

% \fancyhead[RO]{\textbf{\rightmark}}

% \fancyhead[LO]{\textbf{Kai Luo}}
% \fancyhead[LO]{罗凯}

% \fancyhead[LE]{\nouppercase{\rightmark\hfill\leftmark}}
% \fancyhead[RO]{\nouppercase{\leftmark\hfill\rightmark}}

\fancyfoot{} % clear all footer fields
% \fancyfoot[CE,CO]{\rightmark}
% \fancyfoot[LO,CE]{}
% \fancyfoot[CO,RE]{Nanjing University of Science and Technology}
\fancyfoot[C]{\thepage}
\fancyfoot[L]{截止日期\duedate} % clear all footer fields
\fancyfoot[R]{任课教师:罗凯} % clear all footer fields

% \fancyfoot[CE,O]{\thechapter}



\begin{document}
\renewcommand{\labelenumi}{(\arabic{enumi})}
\renewcommand{\labelenumii}{(\arabic{enumi}.\arabic{enumii})}


\begin{problem}
  用爱因斯坦求和约定改写下列方程. 指标取 $1,2,3$，重复指标表示求和，式中不再出现 $\sum$. 叉积和旋度用 $\varepsilon_{ijk}$.
  \begin{enumerate}
    \item 连续性方程 $\partial_t\rho+\nabla\cdot\mathbf{J}=0$;
    \item 泊松方程 $\nabla^2\varphi=-\rho/\varepsilon_0$;
    \item 真空中的麦克斯韦方程组
      \[
        \begin{aligned}
          \nabla\cdot\mathbf{E}&=\rho/\varepsilon_0, &
          \nabla\cdot\mathbf{B}&=0,\\
          \nabla\times\mathbf{E}&=-\partial_t\mathbf{B}, &
          \nabla\times\mathbf{B}
          &=\mu_0\mathbf{J}+\mu_0\varepsilon_0\partial_t\mathbf{E};
        \end{aligned}
      \]
    \item 无源波动方程
      $\nabla^2\mathbf{E}-c^{-2}\partial_{tt}\mathbf{E}=\mathbf{0}$;
    \item 洛伦兹力
      $\mathbf{F}=q\bigl(\mathbf{E}+\mathbf{v}\times\mathbf{B}\bigr)$.
  \end{enumerate}
\end{problem}

\vspace{1.6em}

\begin{problem}
  按爱因斯坦约定去掉求和号，再缩并成更短的形式.
  \begin{enumerate}
    \item $\displaystyle\sum_{j=1}^{3}\delta_{ij}\,\partial_j\varphi$;
    \item $\displaystyle\sum_{i=1}^{3}\sum_{j=1}^{3}\delta_{ij}\,\partial_i A_j$;
    \item $\displaystyle\sum_{j=1}^{3}\sum_{k=1}^{3}\varepsilon_{ijk}\,\delta_{jk}$;
    \item 不可压缩流动
      \[
        \partial_t u_i
        +\sum_{j=1}^{3}u_j\partial_j u_i
        =-\frac{1}{\rho}\sum_{j=1}^{3}\delta_{ij}\,\partial_j p
         +\nu\sum_{j=1}^{3}\partial_j\partial_j u_i,
        \qquad
        \sum_{i=1}^{3}\partial_i u_i=0.
      \]
  \end{enumerate}
\end{problem}
  

% \begin{problem}
% 请利用高斯定理推导出真空中有源静电场方程,请问求解需要初始条件吗？请用文字说明.
% \end{problem}

% \begin{problem}
%   课堂上，我们将$x,y$变量代换成$\xi, \eta$给出了$u_{xy}$的表达式，请推导验证
%   $$
%   u_{x x} 
%   % =\left(u_{\xi \xi} \xi_x^2+u_{\xi \eta} \xi_x \eta_x+u_{\xi} \xi_{x x}\right)+\left(u_{\eta \xi} \eta_x \xi_x+u_{\eta \eta} \eta_x^2+u_\eta \eta_{x x}\right) \\
%   =u_{\xi \xi} \xi_x^2+2 u_{\xi \eta} \xi_x \eta_x+u_{\eta \eta} \eta_x^2+u_{\xi} \xi_{x x}+u_\eta \eta_{x x}.
%   $$.
% \end{problem}
   
% \begin{problem}
%   整数$m$阶贝塞尔方程(Bessel's equation)
%   $$
%     x^2 y^{\prime \prime}+x y^{\prime}+\left(x^2-m^2\right) y = 0, 
%   $$
% 的Laplace变换($m = 0$)为一阶常微分方程
%   $$
%    (p^2 + 1) f'(p) + pf(p) = 0.
%   $$
%    试求解$f(p)$.
% \end{problem}

% \begin{problem}
% 求$$
% (2x + e^{y^2}) dx + (2xy e^{y^2} -2y ) dy = 0
% $$
% 的通解.
% \end{problem}
  



\begin{problem}
  用指标记号证明下列乘积法则. $f$ 是标量场, $\mathbf{A}$、$\mathbf{B}$ 是矢量场, 分量在直角坐标中足够光滑.
  \begin{enumerate}
    \item $\displaystyle
      \nabla\cdot(f\mathbf{A})
      =f\,\nabla\cdot\mathbf{A}+\mathbf{A}\cdot\nabla f$;
    \item $\displaystyle
      \nabla\times(f\mathbf{A})
      =f\,\nabla\times\mathbf{A}+\nabla f\times\mathbf{A}$;
    \item $\displaystyle
      \nabla\cdot(\mathbf{A}\times\mathbf{B})
      =\mathbf{B}\cdot(\nabla\times\mathbf{A})
       -\mathbf{A}\cdot(\nabla\times\mathbf{B})$.
  \end{enumerate}
\end{problem}

\vspace{1.6em}

\begin{problem}
  \begin{enumerate}
    \item 设 $\mathbf{A}=x\,\mathbf{e}_x+y\,\mathbf{e}_y+z\,\mathbf{e}_z$，$V$ 是球 $x^2+y^2+z^2\le R^2$，$S=\partial V$ 取外法向. 分别计算散度定理的两边，验证
      \[
        \oint_S\mathbf{A}\cdot\mathrm{d}\mathbf{S}
        =\int_V(\nabla\cdot\mathbf{A})\,\mathrm{d}V.
      \]
    \item 设 $\mathbf{A}=y\,\mathbf{e}_x+z\,\mathbf{e}_y+x\,\mathbf{e}_z$. $S$ 是正方形 $0\le x\le a$，$0\le y\le a$，$z=0$，法向取 $\mathbf{e}_z$，边界 $L$ 按右手定则定向. 分别计算斯托克斯定理的两边，验证
      \[
        \oint_L\mathbf{A}\cdot\mathrm{d}\boldsymbol{\ell}
        =\int_S(\nabla\times\mathbf{A})\cdot\mathrm{d}\mathbf{S}.
      \]
  \end{enumerate}
\end{problem}

\vspace{1.6em}

\begin{problem}
  已知
  $\mathbf{A}\times(\mathbf{B}\times\mathbf{C})
  =(\mathbf{A}\cdot\mathbf{C})\mathbf{B}-(\mathbf{A}\cdot\mathbf{B})\mathbf{C}$.
  \begin{enumerate}
    \item 证明
      \[
        (\mathbf{A}\times\mathbf{B})\times\mathbf{C}
        =(\mathbf{A}\cdot\mathbf{C})\mathbf{B}
         -(\mathbf{B}\cdot\mathbf{C})\mathbf{A}.
      \]
    \item 证明
      \[
        \mathbf{A}\times(\mathbf{B}\times\mathbf{C})
        +\mathbf{B}\times(\mathbf{C}\times\mathbf{A})
        +\mathbf{C}\times(\mathbf{A}\times\mathbf{B})
        =\mathbf{0}.
      \]
    \item 用指标记号证明拉格朗日恒等式
      \[
        (\mathbf{A}\times\mathbf{B})\cdot(\mathbf{C}\times\mathbf{D})
        =(\mathbf{A}\cdot\mathbf{C})(\mathbf{B}\cdot\mathbf{D})
         -(\mathbf{A}\cdot\mathbf{D})(\mathbf{B}\cdot\mathbf{C}).
      \]
  \end{enumerate}
\end{problem}

\vspace{1.6em}

% \begin{problem}
%   \begin{enumerate}
%     \item 设 $u$ 只依赖于球径 $r$，且 $r>0$. 由球坐标中的拉普拉斯算子证明
%       \[
%         \nabla^2 u
%         =\frac{1}{r^2}\frac{\mathrm{d}}{\mathrm{d}r}
%           \Bigl(r^2\frac{\mathrm{d}u}{\mathrm{d}r}\Bigr).
%       \]
%       对常数 $n$ 计算 $\nabla^2(r^n)$，并求使 $\nabla^2(r^n)=0$（$r>0$）的 $n$.
%     \item 设 $u$ 只依赖于柱径 $\rho$，且 $\rho>0$，与 $\phi$、$z$ 无关. 由柱坐标中的拉普拉斯算子证明
%       \[
%         \nabla^2 u
%         =\frac{1}{\rho}\frac{\mathrm{d}}{\mathrm{d}\rho}
%           \Bigl(\rho\frac{\mathrm{d}u}{\mathrm{d}\rho}\Bigr),
%       \]
%       并计算 $\nabla^2(\rho^n)$.
%   \end{enumerate}
% \end{problem}

\vspace{1.6em}



  
  


% \begin{problem}
% \end{problem}

% \begin{problem}
% \end{problem}

% \begin{problem}
% \end{problem}

\end{document}